\documentclass[11pt]{article}
\usepackage[margin=1in]{geometry}
\usepackage{amsmath,amssymb,bm}
\usepackage{booktabs}
\usepackage{hyperref}

\title{Standalone 0D--2P Relativistic Electron Fokker--Planck Model\\Physical and Numerical Reference}
\author{}
\date{September 2026}

\newcommand{\p}{p}
\newcommand{\x}{\xi}
\newcommand{\g}{\gamma}
\newcommand{\N}{N}
\newcommand{\f}{f}
\newcommand{\Ebar}{\overline{E}}

\begin{document}
\maketitle

\noindent
This document specifies the spatially homogeneous 0D--2P forward relativistic
electron Fokker--Planck model and its GPU finite-volume realization.  The model
includes electric-field acceleration, relativistic small-angle collisions,
synchrotron radiation reaction, and an optional Chiu--Harvey large-angle
secondary-electron source.  Plasma quantities are prescribed and constant for
the present solver.  Future work will exchange state and source information
with a time-dependent bulk plasma model while preserving the conservative
finite-volume and direct-solver contracts defined here.

The numerical reference covers the dimensional and normalized equations, fixed
coefficient construction, conservative cell-content discretization, boundary
conditions, implicit time integration, GPU-resident CSR assembly, cuDSS solves,
and forward-solver qualification.  Its scope ends with forward evolution and
the numerical diagnostics required to qualify that evolution.

Revision v15.4 is reorganized here as a forward-only reference.  The equations
and numerical conventions retained below are the authority for
`forward_fv_solver.py` and `forward_fv_solver.toml`.

%===============================================================================
\section{Physical model}

The present model is a spatially homogeneous, gyrotropic electron kinetic system in
momentum magnitude and pitch angle.  All plasma quantities are prescribed and constant
in time.  Plasma evolution, spatial transport, and kinetic--bulk feedback are outside
the current standalone problem; they are specified as future coupling requirements in
Sec.~\ref{sec:future_bulk_coupling}.  The physical system is defined below in dimensional
form.  No reference normalization used by the numerical solver is introduced until the
final subsection of this section.

\subsection{Dimensional kinetic equation}

The physical model is defined first in dimensional variables.  Let
$F_e(P,\xi,t)$ be the gyrophase-averaged electron distribution function, where $P$ is
the physical momentum magnitude and
\begin{equation}
 \xi \equiv \frac{P_\parallel}{P}\in[-1,1].
\end{equation}
The forward kinetic equation is
\begin{equation}
 \frac{\partial F_e}{\partial t}
 = \mathcal C_{\rm sa}[F_e]
 + \mathcal E[F_e]
 + \mathcal R_{\rm syn}[F_e]
 + \mathcal S_{\rm LA}[F_e].
 \label{eq:dimensional_kinetic}
\end{equation}
Here $\mathcal C_{\rm sa}$ is the relativistic small-angle test-particle
Fokker--Planck collision operator, $\mathcal E$ is acceleration by the parallel
electric field, $\mathcal R_{\rm syn}$ is synchrotron radiation reaction, and
$\mathcal S_{\rm LA}$ is either zero or the Chiu--Harvey large-angle
secondary-electron gain operator defined below.  The same forward operator structure is standard in relativistic runaway-electron
kinetics; see, for example, Ref.~\cite{BraamsKarney1989}.

The dimensional distribution is defined such that its momentum-space integral gives the
represented electron density,
\begin{equation}
 n_F(t)
 = 2\pi\int_0^\infty P^2\,dP\int_{-1}^{1}F_e(P,\xi,t)\,d\xi,
 \label{eq:dimensional_density}
\end{equation}
so that $n_F$ is the represented electron density.  The relativistic Lorentz factor
and speed are
\begin{equation}
 \gamma(P)=\sqrt{1+\frac{P^2}{m_e^2c^2}},
 \qquad
 v(P)=\frac{P}{\gamma m_e},
 \qquad
 \beta(P)=\frac{v}{c}.
 \label{eq:dimensional_kinematics}
\end{equation}

\subsection{Parallel electric-field operator}

For an electron of charge $-e$ in a prescribed uniform parallel electric field
$E_\parallel$, the momentum-space characteristics are
\begin{equation}
 \dot P_E=-eE_\parallel\xi,
 \qquad
 \dot\xi_E=-\frac{eE_\parallel}{P}(1-\xi^2).
 \label{eq:electric_characteristics}
\end{equation}
The corresponding conservative kinetic operator is
\begin{equation}
 \mathcal E[F_e]
 =-\frac{1}{P^2}\frac{\partial}{\partial P}
   \left(P^2\dot P_E F_e\right)
  -\frac{\partial}{\partial\xi}
   \left(\dot\xi_E F_e\right).
 \label{eq:electric_operator_dimensional}
\end{equation}
This is the standard homogeneous relativistic electric-field operator used in
runaway-electron Fokker--Planck models; see
Refs.~\cite{ConnorHastie1975}.

\subsection{Small-angle test-particle collision-model hierarchy}

The standalone kinetic model admits two test-particle approximations for the
small-angle collision kernel.  The production default is the finite-temperature
relativistic model defined first below; the second is the fully-relativistic and
asymptotically matched model historically used in CODE.  These define two alternative
\emph{thermal-log complete-screening base kernels}.  The physical coefficient construction
then applies two common layers that are independent of which base kernel is selected:
(i) energy-dependent Coulomb-logarithm matching of the electron--electron and
electron--ion sectors, and (ii) partial-screening corrections from bound electrons.
Neither layer modifies radial energy diffusion.  Thus changing the base kernel does not
change the Coulomb-logarithm or screening prescription.  Section~2 specifies the same
layered construction in the common solver normalization, without changing the grid, time
integrator, or sparse-solver architecture.

\paragraph{Finite-temperature relativistic test-particle model (current baseline).}
Small-angle collisions are represented by a linearized relativistic test-particle
Fokker--Planck operator with radial drag, radial energy diffusion, and pitch-angle
deflection,
\begin{equation}
 \mathcal C_{\rm sa}[F_e]
 =\frac{1}{P^2}\frac{\partial}{\partial P}
  \left[
   P^2\left(
      C_F^{(P)} F_e
      + C_A^{(P)}\frac{\partial F_e}{\partial P}
   \right)
  \right]
 +\frac{\nu_D}{2}\frac{\partial}{\partial\xi}
  \left[(1-\xi^2)\frac{\partial F_e}{\partial\xi}\right].
 \label{eq:collision_operator_dimensional}
\end{equation}
The dimensional radial coefficients are written in terms of the slowing-down and
parallel diffusion frequencies as
\begin{equation}
 C_F^{(P)}=P\nu_s,
 \qquad
 C_A^{(P)}=\frac{P^2}{2}\nu_\parallel.
 \label{eq:dimensional_collision_coeffs}
\end{equation}
Equation~\eqref{eq:collision_operator_dimensional} is the relativistic test-particle
operator of Braams--Karney/Pike--Rose form.  Its thermal-log complete-screening base
coefficients are written below; the common energy-dependent Coulomb-logarithm and
partial-screening layers are then specified separately following
Refs.~\cite{BraamsKarney1989,PikeRose2014,Hesslow2018Collision}.
The field-particle part of the small-angle collision operator is not retained in the
present model.

The published collision coefficients are most compactly written using the auxiliary
relativistic kinematic variables
\begin{equation}
 u\equiv\frac{P}{m_ec},
 \qquad
 \gamma=\sqrt{1+u^2},
 \qquad
 \Theta\equiv\frac{T_e}{m_ec^2}.
 \label{eq:collision_auxiliary_variables}
\end{equation}
These dimensionless quantities are only arguments used to write the dimensional collision
frequencies compactly; they do not constitute the reference normalization of the kinetic
PDE used by the numerical solver.  Define the physical relativistic collision time
\begin{equation}
 \tau_c
 =\frac{4\pi\epsilon_0^2m_e^2c^3}
 {n_e e^4\ln\Lambda_0},
 \label{eq:physical_collision_time}
\end{equation}
where $n_e$ and $T_e$ are the prescribed free-electron density and temperature and
$\ln\Lambda_0$ is the thermal Coulomb logarithm.  Following the finite-temperature
construction in Ref.~\cite{Hesslow2018Collision}, first define the finite-temperature
\emph{thermal-log} complete-screening base by setting both collisional logarithms equal
to $\ln\Lambda_0$.  It is useful to keep the electron--electron and electron--ion
parts of the deflection frequency separate:
\begin{align}
 \nu_s^{\rm FT,0}
 &=\frac{1}{\tau_c}\frac{\gamma^2}{u^3}
   \frac{u^2}{\gamma^2}\Psi_s(u,\Theta),
 \label{eq:physical_nus}\\
 \nu_\parallel^{\rm FT,0}
 &=\frac{1}{\tau_c}\frac{2\gamma\Theta}{u^3}\Psi_s(u,\Theta),
 \label{eq:physical_nupar}\\
 \nu_{D,ee}^{\rm FT,0}
 &=\frac{1}{\tau_c}\frac{\gamma}{u^3}
   \frac{2u}{\gamma}\Psi_D(u,\Theta),\\
 \nu_{D,ei}^{\rm FT,0}
 &=\frac{1}{\tau_c}\frac{\gamma}{u^3}Z_{\rm eff},\\
 \nu_D^{\rm FT,0}
 &=\nu_{D,ee}^{\rm FT,0}+\nu_{D,ei}^{\rm FT,0}.
 \label{eq:physical_nud}
\end{align}
The finite-temperature relativistic functions are
\begin{align}
 \Psi_s
 &=\frac{
   \gamma^2\Psi_1-\Theta\Psi_0
   +(\Theta\gamma-1)u e^{-\gamma/\Theta}}
  {u^2K_2(1/\Theta)},
 \label{eq:psis_physical}\\
 \Psi_D
 &=\frac{
  (u^2\gamma^2+\Theta^2)\Psi_0
  +\Theta(2u^4-1)\Psi_1
  +\gamma\Theta[1+\Theta(2u^2-1)]u e^{-\gamma/\Theta}}
  {2\gamma u^3K_2(1/\Theta)},
 \label{eq:psid_physical}
\end{align}
with
\begin{equation}
 \Psi_0(u,\Theta)
 =\int_0^u\frac{e^{-\sqrt{1+s^2}/\Theta}}{\sqrt{1+s^2}}\,ds,
 \qquad
 \Psi_1(u,\Theta)
 =\int_0^u e^{-\sqrt{1+s^2}/\Theta}\,ds.
 \label{eq:psi_integrals_physical}
\end{equation}
Here $K_2$ is the modified Bessel function of the second kind.  Numerically stable
evaluation of Eqs.~\eqref{eq:psis_physical}--\eqref{eq:psi_integrals_physical}, including
the low-$u$ limit, belongs to the numerical model in Sec.~2 rather than to the physical
definition of the operator.

The thermal logarithm appearing in Eq.~\eqref{eq:physical_collision_time} is
\begin{equation}
 \ln\Lambda_0
 =14.9-\frac12\ln\left(\frac{n_e}{10^{20}\,{\rm m}^{-3}}\right)
 +\ln\left(\frac{T_e}{1\,{\rm keV}}\right).
 \label{eq:physical_thermal_log}
\end{equation}

\paragraph{Energy-dependent Coulomb-logarithm correction common to both base models.}
The thermal-to-relativistic matching of the Coulomb logarithms is a physical correction
layer, not part of the definition of either small-angle base kernel.  Following Hesslow
et al.~\cite{Hesslow2018Collision}, define
\begin{align}
 \ln\Lambda_{ee}(u)
 &=\ln\Lambda_0+\frac{1}{k}
 \ln\left[
  1+\left(\frac{2(\gamma-1)}{u_{Te}^2}\right)^{k/2}
 \right],\\
 \ln\Lambda_{ei}(u)
 &=\ln\Lambda_0+\frac{1}{k}
 \ln\left[1+\left(\frac{2u}{u_{Te}}\right)^k\right],
 \qquad
 u_{Te}=\sqrt{2\Theta},
 \qquad k=5,
 \label{eq:physical_matched_logs}
\end{align}
and the ratios
\begin{equation}
 R_{ee}(u)=\frac{\ln\Lambda_{ee}(u)}{\ln\Lambda_0},
 \qquad
 R_{ei}(u)=\frac{\ln\Lambda_{ei}(u)}{\ln\Lambda_0}.
 \label{eq:physical_log_ratios}
\end{equation}
For either base $b\in\{\mathrm{FT},\mathrm{match}\}$, decompose its thermal-log
complete-screening deflection frequency into electron--electron and electron--ion parts,
$\nu_D^{b,0}=\nu_{D,ee}^{b,0}+\nu_{D,ei}^{b,0}$.  The common Coulomb-logarithm layer is
then
\begin{align}
 \nu_s^{b,\rm CL}&=R_{ee}\,\nu_s^{b,0},\\
 \nu_\parallel^{b,\rm CL}&=\nu_\parallel^{b,0},\\
 \nu_D^{b,\rm CL}
 &=R_{ee}\,\nu_{D,ee}^{b,0}
  +R_{ei}\,\nu_{D,ei}^{b,0}.
 \label{eq:physical_common_log_overlay}
\end{align}
Thus the electron--electron logarithm corrects collisional drag and the
 electron--electron part of pitch-angle deflection, the electron--ion logarithm corrects
only the electron--ion deflection sector, and radial energy diffusion is unchanged.
For the finite-temperature base the split is explicit in
Eqs.~\eqref{eq:physical_nus}--\eqref{eq:physical_nud}.  For the matched base, the
$Z_{\rm eff}$ part of the deflection coefficient is the electron--ion sector and the
remaining deflection terms are the electron--electron sector, as made explicit in the
solver-normalized formulas in Sec.~2.

\paragraph{Partial-screening correction common to both base models.}
For a partially ionized ion species $a$ with atomic number $Z_a$, ion charge
$Z_{0a}$, density $n_a$, and number of bound electrons $N_a=Z_a-Z_{0a}$, define
\cite{Hesslow2017Screening,Hesslow2018Collision}
\begin{align}
 h(u)
 &=\sum_a\frac{n_a}{n_e}N_a
 \left\{
  \frac15\ln\left[
   1+\left(\frac{u\sqrt{\gamma-1}}{I_a/(m_ec^2)}\right)^5
  \right]-\beta^2
 \right\},
 \label{eq:physical_h}\\
 g(u)
 &=\sum_a\frac{n_a}{n_e}
 \left\{
  \frac23(Z_a^2-Z_{0a}^2)
       \ln\left[1+(u\bar a_a)^{3/2}\right]
  -\frac23N_a^2
       \frac{(u\bar a_a)^{3/2}}{1+(u\bar a_a)^{3/2}}
 \right\}.
 \label{eq:physical_g}
\end{align}
Here $I_a$ is the mean excitation energy and $\bar a_a$ is the dimensionless screening
length defined in Ref.~\cite{Hesslow2018Collision}.  The corresponding physical
frequency corrections are
\begin{align}
 \Delta\nu_s^{\rm PS}
 &=\frac{1}{\tau_c}\frac{\gamma^2}{u^3}
   \frac{h(u)}{\ln\Lambda_0},\\
 \Delta\nu_\parallel^{\rm PS}&=0,\\
 \Delta\nu_D^{\rm PS}
 &=\frac{1}{\tau_c}\frac{\gamma}{u^3}
   \frac{g(u)}{\ln\Lambda_0}.
 \label{eq:physical_partial_screening_overlay}
\end{align}
For either base kernel, the partial-screening correction is applied after the common
energy-dependent Coulomb-logarithm layer.  Equivalently,
$\Delta C_F^{(P),\rm PS}=P\Delta\nu_s^{\rm PS}$,
$\Delta C_A^{(P),\rm PS}=0$, and the deflection coefficient receives
$\Delta\nu_D^{\rm PS}$.  The final physical frequencies are therefore
\begin{align}
 \nu_s^b&=\nu_s^{b,\rm CL}+\Delta\nu_s^{\rm PS},\\
 \nu_\parallel^b&=\nu_\parallel^{b,\rm CL}=\nu_\parallel^{b,0},\\
 \nu_D^b&=\nu_D^{b,\rm CL}+\Delta\nu_D^{\rm PS},
 \qquad b\in\{\mathrm{FT},\mathrm{match}\}.
 \label{eq:physical_final_collision_layers}
\end{align}
Thus neither the energy-dependent Coulomb-logarithm layer nor the partial-screening layer
selects or blends the base operator, and neither modifies radial energy diffusion.

The effective charge is
\begin{equation}
 Z_{\rm eff}=\frac{1}{n_e}\sum_a n_a Z_{0a}^2.
\end{equation}
In the complete-screening limit $N_a=0$ for every species, the screening corrections
vanish identically, $h=g=0$, while the common energy-dependent Coulomb-logarithm layer
remains active.  Setting additionally $R_{ee}=R_{ei}=1$ recovers the corresponding
thermal-log base kernel exactly.


\paragraph{Fully-relativistic and asymptotically matched test-particle model.}
The second selectable thermal-log complete-screening base collision kernel is the fully-relativistic,
asymptotically matched test-particle operator.  Stahl et al. describe its construction as
an asymptotic matching of the usual non-relativistic test-particle Fokker--Planck
operator to the fully relativistic high-energy test-particle limit
\cite{Stahl2016}.  The resulting momentum-space coefficients retain relativistic
kinematics throughout the represented tail while assuming a non-relativistic thermal
bulk.  This is historically the test-particle operator used in CODE, but CODE is
provenance rather than the model name.  It is not the finite-temperature
Braams--Karney/Pike--Rose operator above, and it does not contain the field-particle part
of the linearized collision operator.

For provenance and comparison with Stahl et al., it is useful to state this model in
the normalized variables used in their CODE implementation.  Choose a reference thermal speed
\begin{equation}
 \widetilde v_e=\sqrt{\frac{2\widetilde T_e}{m_e}},\qquad
 \delta\equiv\frac{\widetilde v_e}{c},
\end{equation}
and define
\begin{equation}
 y\equiv\frac{P}{m_e\widetilde v_e}=\frac{u}{\delta},\qquad
 x\equiv\frac{y}{\gamma}=\frac{v}{\widetilde v_e},\qquad
 \overline v_e\equiv\frac{v_e}{\widetilde v_e},\qquad
 v_e=\sqrt{\frac{2T_e}{m_e}}.
 \label{eq:matched_variables}
\end{equation}
With time normalized to a reference electron--electron collision frequency
$\widetilde\nu_{ee}$, the fully-relativistic asymptotically matched test-particle operator is
\begin{equation}
 \boxed{
 \widehat{\mathcal C}_{\rm match}^{\rm tp}[G]
 =c_C\overline v_e^{\,3}y^{-2}
 \left\{
 \frac{\partial}{\partial y}
 \left[
  y^2\Psi
  \left(
   \frac{1}{x}\frac{\partial}{\partial y}
   +\frac{2}{\overline v_e^{\,2}}
  \right)G
 \right]
 +\frac{c_\xi}{2x}\frac{\partial}{\partial\xi}
 \left[(1-\xi^2)\frac{\partial G}{\partial\xi}\right]
 \right\},
 }
 \label{eq:matched_tp_operator}
\end{equation}
where a bar denotes normalization to the reference plasma state, in particular
$\overline\nu_{ee}=\nu_{ee}/\widetilde\nu_{ee}$, and
\begin{align}
 c_C&=\frac{3\sqrt{\pi}}{4}\,\overline\nu_{ee},\\
 c_\xi&=Z_{\rm eff}+\Phi-\Psi
       +\frac{\overline v_e^{\,2}\delta^4x^2}{2},
 \label{eq:matched_coefficients}
\end{align}
with
\begin{equation}
 \Phi=\operatorname{erf}\!\left(\frac{x}{\overline v_e}\right),\qquad
 \Psi=\frac{\overline v_e^{\,2}}{2x^2}
 \left[
  \Phi-\frac{x}{\overline v_e}
  \frac{d\Phi}{d(x/\overline v_e)}
 \right].
 \label{eq:matched_chandrasekhar}
\end{equation}
For a constant plasma chosen equal to the reference state,
$\overline v_e=\overline\nu_{ee}=1$.  The assumption $\delta\ll1$ is the
non-relativistic-bulk ordering of this model.  The role of
Eq.~\eqref{eq:matched_tp_operator} in the model hierarchy is specifically to
recover the standard non-relativistic test-particle coefficients at thermal and
suprathermal energies while matching to the relativistic test-particle asymptote at
large momentum \cite{Stahl2016}.  This provides a useful fully-relativistic asymptotically matched configuration for
benchmarks without replacing the finite-temperature relativistic baseline above.  For this matched base, the slowing-down coefficient is electron--electron, the
$Z_{\rm eff}$ contribution to pitch-angle deflection is electron--ion, and the remaining
pitch-angle-deflection terms are electron--electron.  Consequently the same common
Coulomb-logarithm layer in Eq.~\eqref{eq:physical_common_log_overlay} applies to this
base without introducing a matched-model-specific logarithm prescription.  When bound
electrons are present, the same Eq.~\eqref{eq:physical_partial_screening_overlay}
correction is then added; no separate screening prescription is introduced either.

\subsection{Synchrotron radiation-reaction operator}

The gyro-averaged synchrotron radiation-reaction force is represented by the
momentum-space characteristics
\begin{equation}
 \dot P_{\rm syn}
 =-\frac{\gamma P}{\tau_s}(1-\xi^2),
 \qquad
 \dot\xi_{\rm syn}
 =\frac{\xi}{\gamma\tau_s}(1-\xi^2),
 \label{eq:synch_characteristics}
\end{equation}
with
\begin{equation}
 \tau_s=\frac{6\pi\epsilon_0m_e^3c^3}{e^4B^2}.
 \label{eq:synch_time}
\end{equation}
The conservative operator is
\begin{equation}
 \mathcal R_{\rm syn}[F_e]
 =-\frac{1}{P^2}\frac{\partial}{\partial P}
   \left(P^2\dot P_{\rm syn}F_e\right)
  -\frac{\partial}{\partial\xi}
   \left(\dot\xi_{\rm syn}F_e\right).
 \label{eq:synch_operator_dimensional}
\end{equation}
This is the standard gyro-averaged synchrotron reaction used in homogeneous
runaway-electron kinetic models; see Refs.~\cite{Andersson2001}.

\subsection{Chiu--Harvey large-angle secondary-electron gain}

The production model uses a single optional large-angle secondary-electron gain model:
the finite-primary-energy, aligned-primary Chiu--Harvey source
\cite{Chiu1998,Harvey2000}.  The source is gain only: it adds the represented
secondary electron but does not include the compensating recoil/loss of the primary.
The source may be disabled by setting $\mathcal S_{\rm LA}=0$.

Introduce the dimensionless momenta
\begin{equation}
 u\equiv\frac{P}{m_ec},\qquad
 u'\equiv\frac{P'}{m_ec},\qquad
 \gamma=\sqrt{1+u^2},\qquad
 \gamma'=\sqrt{1+u'^2},
 \label{eq:ch_auxiliary_momenta}
\end{equation}
and the rescaled distribution
\begin{equation}
 \widehat F_e(u,\xi,t)\equiv(m_ec)^3F_e(P,\xi,t).
 \label{eq:ch_hat_definition}
\end{equation}
This is only a change of momentum coordinate; it is not the reference normalization of
Sec.~\ref{sec:normalization}.  Define the pitch-integrated primary spectrum
\begin{equation}
 \boxed{
 \mathcal F(u',t)=\int_{-1}^{1}\widehat F_e(u',\xi',t)\,d\xi'.
 }
 \label{eq:ch_primary_spectrum}
\end{equation}
The Chiu--Harvey approximation retains this evolving primary energy spectrum but uses a
field-aligned primary direction, $\xi'=-1$, in the relativistic electron--electron
knock-on kinematics.

The normalized M{\o}ller differential cross section underlying the Chiu--Harvey source is
\cite{Moller1932,Chiu1998,Harvey2000}
\begin{equation}
 \frac{d\overline\sigma_{\rm ee}(\gamma',\gamma)}{du}
 =2\pi
 \frac{\beta\,\gamma'^2}{(\gamma'-1)^3(\gamma'+1)}
 \left[
  x^2-3x+\left(\frac{\gamma'-1}{\gamma'}\right)^2(1+x)
 \right],
 \label{eq:ch_knockon_cross_section}
\end{equation}
where $\overline\sigma_{\rm ee}=\sigma_{\rm ee}/r_e^2$,
\begin{equation}
 \beta=\frac{u}{\gamma},\qquad
 x=\frac{(\gamma'-1)^2}{(\gamma-1)(\gamma'-\gamma)},\qquad
 r_e=\frac{e^2}{4\pi\epsilon_0m_ec^2}.
 \label{eq:ch_knockon_cross_section_aux}
\end{equation}
For the aligned primary, the secondary pitch is
\begin{equation}
 \xi_{\rm CH}(\gamma',\gamma)
 =-\left[
  \frac{(\gamma'+1)(\gamma-1)}{(\gamma'-1)(\gamma+1)}
 \right]^{1/2}.
 \label{eq:ch_secondary_pitch}
\end{equation}
The corresponding gain operator is
\begin{equation}
 \widehat S_{\rm CH}(u,\xi,t)
 =\frac{n_t c r_e^2}{u^2}
  \int_0^\infty du'\,u'^2\beta'
  \frac{d\overline\sigma_{\rm ee}(\gamma',\gamma)}{du}
  \delta\!\left[\xi-\xi_{\rm CH}(\gamma',\gamma)\right]
  \mathcal F(u',t),
 \label{eq:ch_integral_source}
\end{equation}
with $\beta'=u'/\gamma'$.  Thus no prescribed analytic primary spectrum is introduced:
$\mathcal F$ is obtained directly from the evolving 2P solution.

Equation~\eqref{eq:ch_integral_source} can be reduced analytically over $u'$.  The
required primary root is
\begin{equation}
 u'_0
 =-\frac{2u\xi}{1+\xi^2-\gamma(1-\xi^2)}
 =\frac{2u|\xi|}{1+\xi^2-\gamma(1-\xi^2)},
 \qquad \xi<0,
 \label{eq:ch_primary_root}
\end{equation}
with Jacobian
\begin{equation}
 \left|
 \frac{\partial(\xi-\xi_{\rm CH})}{\partial u'}
 \right|_{u'=u'_0}
 =\frac{|\xi|}{u'_0\gamma'_0},
 \qquad \gamma'_0=\sqrt{1+u_0'^2}.
 \label{eq:ch_delta_jacobian}
\end{equation}
Using $\beta'_0=u'_0/\gamma'_0$ gives
\begin{equation}
 \boxed{
 \widehat S_{\rm CH}(u,\xi,t)
 =\frac{n_t c r_e^2}{u^2}
  \frac{d\overline\sigma_{\rm ee}(u'_0,u)}{du}
  \frac{u_0'^4}{|\xi|}\,
  \mathcal F(u'_0,t),
 }
 \label{eq:ch_reduced_source}
\end{equation}
when the root satisfies the kinematic and domain constraints, and zero otherwise.

A fixed large-angle cutoff $P_m$ avoids double counting scattering already represented
by the small-angle Fokker--Planck operator.  With
\begin{equation}
 \epsilon_m=\sqrt{1+(P_m/m_ec)^2}-1,
\end{equation}
the retained secondary kinetic energy $\epsilon_s=\gamma-1$ for a primary with
$\epsilon_p=\gamma'-1$ satisfies
\begin{equation}
 \epsilon_s\in\left[\epsilon_m,\frac{\epsilon_p}{2}\right].
 \label{eq:ch_energy_range}
\end{equation}
The source is also zero if the root in Eq.~\eqref{eq:ch_primary_root} lies outside the
represented primary radial domain.

The dimensional large-angle term is therefore simply
\begin{equation}
 \boxed{
 \mathcal S_{\rm LA}[F_e](P,\xi,t)
 \in\left\{0,\;\frac{1}{(m_ec)^3}\widehat S_{\rm CH}(u,\xi,t)\right\}.
 }
 \label{eq:large_angle_physical_choices}
\end{equation}
The target density $n_t$ is explicit and denotes the population treated as cold
stationary target electrons for the knock-on source.  For the partially ionized
application of Ref.~, $n_t$ is the total free-plus-bound electron
density; in the standalone solver it is an explicit prescribed physical input.

\subsection{Prescribed plasma state, initial condition, and moments}

For the present constant-parameter problem, the prescribed quantities
\begin{equation}
 \left\{T_e,n_e,E_\parallel,B,n_a,Z_a,Z_{0a},I_a,\bar a_a,n_t,P_m\right\}
\end{equation}
are constant in time.  Consequently every coefficient in
Eq.~\eqref{eq:dimensional_kinetic} is time independent.

The initial condition is an arbitrary non-negative gyrotropic distribution
\begin{equation}
 F_e(P,\xi,0)=F_{e0}(P,\xi).
 \label{eq:dimensional_initial_condition}
\end{equation}
Useful physical moments include the represented density in
Eq.~\eqref{eq:dimensional_density} and the parallel current density
\begin{equation}
 j_\parallel(t)
 =-e\,2\pi\int_0^\infty P^2\,dP\int_{-1}^{1}
  v(P)\,\xi\,F_e(P,\xi,t)\,d\xi.
 \label{eq:dimensional_current}
\end{equation}
Additional moments may be formed from the same distribution without modifying the
kinetic equation.

\subsection{Normalization used by the numerical model}
\label{sec:normalization}

Only after the dimensional physical model has been defined do we introduce the
normalization used in Sec.~2.  The local flux form below is common to both specified
test-particle small-angle configurations; Sec.~2 supplies the model-specific values of
$C_F$, $C_A$, and $\nu_D$.  The large-angle term is selected independently as either none or Chiu--Harvey.  Choose a fixed reference density $n_{\rm ref}$
and reference Coulomb logarithm $\ln\Lambda_{\rm ref}$ and define
\begin{equation}
 \tau_{\rm ref}
 =\frac{4\pi\epsilon_0^2m_e^2c^3}
 {n_{\rm ref}e^4\ln\Lambda_{\rm ref}},
 \qquad
 E_{\rm ref}=\frac{m_ec}{e\tau_{\rm ref}}.
 \label{eq:reference_scales}
\end{equation}
The dimensionless numerical variables are
\begin{equation}
 p=\frac{P}{m_ec},
 \qquad
 \tau=\frac{t}{\tau_{\rm ref}},
 \qquad
 f(p,\xi,\tau)=\frac{(m_ec)^3}{n_{\rm ref}}F_e(P,\xi,t),
 \qquad
 \bar E=\frac{E_\parallel}{E_{\rm ref}}.
 \label{eq:numerical_normalization}
\end{equation}
The normalization is fixed for the calculation and does not alter the dimensional
physics.  In particular,
\begin{equation}
 \frac{n_F}{n_{\rm ref}}
 =2\pi\int_0^\infty p^2\,dp\int_{-1}^{1}f(p,\xi,\tau)\,d\xi.
 \label{eq:normalized_density}
\end{equation}
Define normalized collision coefficients
\begin{equation}
 C_F=p\tau_{\rm ref}\nu_s,
 \qquad
 C_A=\frac{p^2}{2}\tau_{\rm ref}\nu_\parallel,
 \qquad
 \nu_D^{\ast}=\tau_{\rm ref}\nu_D,
 \qquad
 \alpha=\frac{\tau_{\rm ref}}{\tau_s}.
 \label{eq:normalized_coefficients}
\end{equation}
In the remainder of the document the star is dropped and $\nu_D$ denotes the
normalized deflection frequency.

With these definitions the normalized forward equation used by the numerical model is
\begin{equation}
 \frac{\partial f}{\partial\tau}
 =-\frac{1}{p^2}\frac{\partial}{\partial p}\left(p^2J_p\right)
  -\frac{\partial J_\xi}{\partial\xi}
  +S_{\rm LA}[f],
 \label{eq:normalized_forward}
\end{equation}
with local fluxes
\begin{align}
 J_p
 &=A_p f-C_A\frac{\partial f}{\partial p},
 \label{eq:normalized_Jp}\\
 J_\xi
 &=A_\xi f
   -\frac{\nu_D}{2}(1-\xi^2)\frac{\partial f}{\partial\xi},
 \label{eq:normalized_Jxi}
\end{align}
and
\begin{align}
 A_p
 &=-\bar E\,\xi
   -\alpha\gamma p(1-\xi^2)
   -C_F,
 \label{eq:normalized_Ap}\\
 A_\xi
 &=(1-\xi^2)
   \left(-\frac{\bar E}{p}+\alpha\frac{\xi}{\gamma}\right).
 \label{eq:normalized_Axi}
\end{align}
The normalized knock-on event-rate scale follows directly from
Eq.~\eqref{eq:ch_integral_source}.  Factoring the $2\pi$ contained in the normalized
differential cross section in Eq.~\eqref{eq:ch_knockon_cross_section},
\begin{equation}
 \tau_{\rm ref}\,2\pi n_t c r_e^2
 =\frac{n_t/n_{\rm ref}}{2\ln\Lambda_{\rm ref}}.
 \label{eq:normalized_knockon_scale}
\end{equation}
This normalization is used by the Chiu--Harvey source.  For constant prescribed plasma
parameters, all selected local coefficients and the Chiu--Harvey map are constant in
time.

\section{Numerical model}
%===============================================================================

The numerical problem is the constant-coefficient normalized forward equation
\eqref{eq:normalized_forward}.  The small-angle configuration selects either the
finite-temperature relativistic or fully-relativistic asymptotically matched thermal-log
base coefficient set.  Independently of that choice, the common energy-dependent
Coulomb-logarithm layer is applied to the electron--electron and electron--ion sectors,
and bound electrons activate the common partial-screening drag/deflection layer defined
in Sec.~1.  Fully stripped species therefore remove only the partial-screening layer; the
energy-dependent Coulomb logarithms remain part of the physical collision coefficients.
The large-angle gain is either disabled or Chiu--Harvey.  Every configuration uses the same physical
finite-volume state in the phase-space measure $2\pi p^2\,dp\,d\xi$ and is advanced as
part of one fully implicit linear kinetic operator.  For fixed physical parameters,
configuration, mesh, and timestep, the semidiscrete operator is linear and time
independent.

\subsection{Finite-volume grid, representative points, and discrete state}

The computational domain is
\begin{equation}
 p_{\min}\le p\le p_{\max},\qquad -1\le\xi\le1,
\end{equation}
and is discretized with a uniform tensor-product mesh.  The production default is
$p_{\min}=0$; a shifted $p_{\min}>0$ domain is available for suprathermal benchmark
configurations.  The radial and pitch spacings are
\begin{equation}
 \Delta p=\frac{p_{\max}-p_{\min}}{N_p},\qquad
 \Delta\xi=\frac{2}{N_\xi}.
 \label{eq:uniform_mesh_spacing}
\end{equation}
The radial faces and cell centers are
\begin{equation}
 p_{i+1/2}=p_{\min}+i\,\Delta p,\qquad
 p_i=p_{\min}+\left(i-\frac12\right)\Delta p,
 \qquad i=1,\ldots,N_p,
 \label{eq:uniform_p_grid}
\end{equation}
and the pitch faces and cell centers are
\begin{equation}
 \xi_{j+1/2}=-1+j\,\Delta\xi,\qquad
 \xi_j=-1+\left(j-\frac12\right)\Delta\xi,
 \qquad j=1,\ldots,N_\xi.
 \label{eq:uniform_xi_grid}
\end{equation}
No radial or pitch stretching is used in the baseline solver.  Resolution is obtained
by increasing $N_p$ and $N_\xi$.  In particular, because the radial drift can remain
strong at relativistic momentum and first-order upwinding introduces numerical diffusion
proportional to the local mesh spacing, $\Delta p$ must be made sufficiently small by
using a large radial mesh.  The use of a GPU is intended to make such uniformly fine
meshes practical.

The primary finite-volume unknown is the exact normalized particle content of cell
$(i,j)$,
\begin{equation}
 N_{ij}(\tau)
 =2\pi\int_{p_{i-1/2}}^{p_{i+1/2}}p^2\,dp
       \int_{\xi_{j-1/2}}^{\xi_{j+1/2}}f(p,\xi,\tau)\,d\xi.
 \label{eq:cellcontent}
\end{equation}
The exact cell phase volume and piecewise-constant reconstruction are
\begin{align}
 V_{ij}
 &=\frac{2\pi}{3}
   \left(p_{i+1/2}^3-p_{i-1/2}^3\right)
   \Delta\xi,
 \label{eq:cellvolume}\\
 f_{ij}&=\frac{N_{ij}}{V_{ij}}.
 \label{eq:cell_reconstruction}
\end{align}
When $p_{\min}=0$, the first radial cell touches the origin and no collision coefficient
is required at $p=0$ itself because the inner flux is imposed by regularity.  For a
shifted domain, the first cell instead terminates on the prescribed inner boundary.

The continuous initial condition is projected into this same finite-volume measure,
\begin{equation}
 N_{ij}^{0}
 =2\pi\int_{p_{i-1/2}}^{p_{i+1/2}}p^2\,dp
       \int_{\xi_{j-1/2}}^{\xi_{j+1/2}}f_0(p,\xi)\,d\xi.
 \label{eq:initial_projection}
\end{equation}
The projection is evaluated analytically when possible and otherwise by converged
Gauss--Legendre quadrature; point sampling of $f_0$ at the cell center is not used as
a substitute for Eq.~\eqref{eq:initial_projection}.

\subsection{Numerical evaluation of the fixed physical coefficients}

All physical coefficients are time independent for the present problem and are therefore
evaluated once before the kinetic operator is assembled.  Both implemented small-angle
models are reduced to the same three normalized coefficient arrays $C_F(p)$, $C_A(p)$,
and $\nu_D(p)$ consumed by the finite-volume fluxes below.  For either base model the
coefficient pipeline is
\begin{equation}
 \begin{aligned}
 \text{thermal-log base}
 &\;\longrightarrow\;\text{common energy-dependent Coulomb logs}\\
 &\;\longrightarrow\;\text{common partial screening}
 \;\longrightarrow\;\{C_F,C_A,\nu_D\}.
 \end{aligned}
 \label{eq:collision_layer_pipeline}
\end{equation}
Selecting a collision model therefore changes only the base coefficient generation, not
the two common correction layers or the finite-volume stencil.  Define once for both
bases
\begin{equation}
 \chi_c=\frac{n_e}{n_{\rm ref}}
        \frac{\ln\Lambda_0}{\ln\Lambda_{\rm ref}}.
 \label{eq:common_collision_scaling}
\end{equation}

\paragraph{Finite-temperature relativistic coefficient evaluation.}
The finite-temperature collision functions in
Eqs.~\eqref{eq:psis_physical}--\eqref{eq:psi_integrals_physical} must be evaluated in a
form that remains accurate both at low temperature and as $p\rightarrow0$.  The
numerical treatment below is an alternative evaluation of the same functions defined in
Sec.~1; it does not introduce a different collision model.

\subparagraph{Exponentially scaled rapidity representation.}
With
\begin{equation}
 p=\sinh y,\qquad \gamma=\cosh y,
 \label{eq:rapidity_substitution}
\end{equation}
define
\begin{align}
 \widetilde\Psi_0(p,\Theta)
 &\equiv e^{1/\Theta}\Psi_0(p,\Theta)
 =\int_0^{\operatorname{asinh}p}
   e^{-[\cosh y-1]/\Theta}\,dy,
 \label{eq:psi0_scaled}\\
 \widetilde\Psi_1(p,\Theta)
 &\equiv e^{1/\Theta}\Psi_1(p,\Theta)
 =\int_0^{\operatorname{asinh}p}
   e^{-[\cosh y-1]/\Theta}\cosh y\,dy,
 \label{eq:psi1_scaled}\\
 \widetilde K_2(\Theta)
 &\equiv e^{1/\Theta}K_2(1/\Theta)
 =\int_0^\infty
   e^{-[\cosh y-1]/\Theta}\cosh(2y)\,dy.
 \label{eq:k2_scaled}
\end{align}
The scaled collision functions are therefore evaluated directly as
\begin{align}
 \Psi_s^{\rm dir}(p,\Theta)
 &=\frac{
   \gamma^2\widetilde\Psi_1-\Theta\widetilde\Psi_0
   +(\Theta\gamma-1)p e^{-(\gamma-1)/\Theta}}
  {p^2\widetilde K_2},
 \label{eq:psis_scaled_direct}\\
 \Psi_D^{\rm dir}(p,\Theta)
 &=\frac{
  (p^2\gamma^2+\Theta^2)\widetilde\Psi_0
  +\Theta(2p^4-1)\widetilde\Psi_1
  +\gamma\Theta[1+\Theta(2p^2-1)]p e^{-(\gamma-1)/\Theta}}
  {2\gamma p^3\widetilde K_2}.
 \label{eq:psid_scaled_direct}
\end{align}
Thus neither $e^{-1/\Theta}$ nor $K_2(1/\Theta)$ is formed separately.  The finite
integrals in Eqs.~\eqref{eq:psi0_scaled} and \eqref{eq:psi1_scaled} are evaluated by
Gauss--Legendre quadrature in $y$; the collision quadrature order $q_{\rm coll}$ is a
numerical convergence parameter.  The same converged value of $\widetilde K_2$ is used
in both the direct and low-$p$ branches.

\subparagraph{Low-$p$ series.}
Although Eqs.~\eqref{eq:psis_scaled_direct} and \eqref{eq:psid_scaled_direct} are free of
thermal underflow, their numerators contain cancellations as $p\rightarrow0$.  The
functions are therefore evaluated near the origin by their Taylor series obtained
directly from the same equations.  Write
\begin{align}
 \Psi_s^{\rm ser}(p,\Theta)
 &=\frac{p}{\widetilde K_2}
   \sum_{m=0}^{6}a_{2m+1}(\Theta)p^{2m},
 \label{eq:psis_lowp_series}\\
 \Psi_D^{\rm ser}(p,\Theta)
 &=\frac{1}{\widetilde K_2}
   \sum_{m=0}^{6}b_{2m}(\Theta)p^{2m}.
 \label{eq:psid_lowp_series}
\end{align}
The retained coefficients, corresponding to terms through $p^{13}$ in $\Psi_s$ and
$p^{12}$ in $\Psi_D$, are
\begin{align}
 a_1&=\frac{2\Theta^2+2\Theta+1}{3\Theta},\\
 a_3&=-\frac{6\Theta^3+6\Theta^2+5\Theta+3}{30\Theta^2},\\
 a_5&=\frac{30\Theta^4+30\Theta^3+27\Theta^2+17\Theta+5}{280\Theta^3},\\
 a_7&=-\frac{210\Theta^5+210\Theta^4+195\Theta^3+125\Theta^2+44\Theta+7}
 {3024\Theta^4},\\
 a_9&=\frac{1890\Theta^6+1890\Theta^5+1785\Theta^4+1155\Theta^3
 +435\Theta^2+92\Theta+9}{38016\Theta^5},\\
 a_{11}&=-\frac{20790\Theta^7+20790\Theta^6+19845\Theta^5+12915\Theta^4
 +5040\Theta^3+1197\Theta^2+167\Theta+11}{549120\Theta^6},\\
 a_{13}&=\frac{270270\Theta^8+270270\Theta^7+259875\Theta^6+169785\Theta^5
 +67725\Theta^4+17136\Theta^3+2786\Theta^2+275\Theta+13}
 {8985600\Theta^7},
 \label{eq:psis_lowp_coeffs}
\end{align}
and
\begin{align}
 b_0&=\frac{8\Theta^2+5\Theta+4}{12},\\
 b_2&=\frac{64\Theta^3+79\Theta^2+30\Theta-8}{240\Theta},\\
 b_4&=-\frac{768\Theta^4+873\Theta^3+274\Theta^2+53\Theta-12}
 {3360\Theta^2},\\
 b_6&=\frac{49152\Theta^5+53877\Theta^4+15594\Theta^3+3830\Theta^2
 +572\Theta-80}{241920\Theta^3},\\
 b_8&=-\frac{983040\Theta^6+1055805\Theta^5+292986\Theta^4+75031\Theta^3
 +16876\Theta^2+1545\Theta-140}{5322240\Theta^4},\\
 b_{10}&=\frac{47185920\Theta^7+50023755\Theta^6+13539990\Theta^5
 +3487020\Theta^4+866196\Theta^3+136817\Theta^2+8050\Theta-504}
 {276756480\Theta^5},\\
 b_{12}&=-\frac{
 \begin{aligned}
  &440401920\Theta^8+462699195\Theta^7+123237750\Theta^6
   +31617235\Theta^5+8090500\Theta^4\\
  &\qquad{}+1515314\Theta^3+167034\Theta^2+6755\Theta-308
 \end{aligned}
 }{2767564800\Theta^6}.
 \label{eq:psid_lowp_coeffs}
\end{align}
For numerical evaluation the two polynomials are evaluated in $z=p^2$ using Horner's
rule.

The baseline branch switch is fixed as
\begin{equation}
 p_{\rm sw}(\Theta)=0.1\min\!\left(1,\sqrt{\Theta}\right),
 \label{eq:lowp_switch}
\end{equation}
with
\begin{equation}
 (\Psi_s,\Psi_D)=
 \begin{cases}
  (\Psi_s^{\rm ser},\Psi_D^{\rm ser}), & p\le p_{\rm sw}(\Theta),\\
  (\Psi_s^{\rm dir},\Psi_D^{\rm dir}), & p>p_{\rm sw}(\Theta).
 \end{cases}
 \label{eq:collision_function_branch}
\end{equation}
This switch places the transition in a region where both $p$ and the thermal-scaled
momentum $p/\sqrt{\Theta}$ are small.  It is fixed for the finite-temperature configuration rather than
being treated as a user-adjustable modeling parameter.  The implementation must recover
\begin{equation}
 \Psi_s=O(p),\qquad \Psi_D=O(1),\qquad p\rightarrow0,
 \label{eq:low_p_limits}
\end{equation}
and the direct and series branches must agree in the neighborhood of
Eq.~\eqref{eq:lowp_switch} to the accuracy required of the coefficient evaluation.

The corresponding normalized finite-temperature thermal-log base coefficients are
\begin{align}
 C_F^{\rm FT,0}(p)
 &=\chi_c\,\Psi_s(p,\Theta),
 \label{eq:ft_numeric_cf_base}\\
 C_A^{\rm FT,0}(p)
 &=\chi_c\,\frac{\gamma\Theta}{p}\Psi_s(p,\Theta),
 \label{eq:ft_numeric_ca_base}\\
 \nu_{D,ee}^{\rm FT,0}(p)
 &=\chi_c\,\frac{2\Psi_D(p,\Theta)}{p^2},\\
 \nu_{D,ei}^{\rm FT,0}(p)
 &=\chi_c\,\frac{\gamma Z_{\rm eff}}{p^3},\\
 \nu_D^{\rm FT,0}(p)
 &=\nu_{D,ee}^{\rm FT,0}(p)+\nu_{D,ei}^{\rm FT,0}(p).
 \label{eq:ft_numeric_nud_base}
\end{align}
These arrays contain the base-kernel dependence only; the common Coulomb-logarithm and
partial-screening layers are applied below.

\paragraph{Fully-relativistic asymptotically matched coefficient evaluation.}
For this configuration the reference plasma is chosen equal to the prescribed
constant plasma state, so that $\overline v_e=\overline\nu_{ee}=1$ in
Eqs.~\eqref{eq:matched_variables}--\eqref{eq:matched_coefficients}.  Define
\begin{equation}
 \delta=\sqrt{\frac{2T_e}{m_ec^2}},\qquad
 x(p)=\frac{p}{\delta\gamma}.
 \label{eq:matched_numeric_scalings}
\end{equation}
Transforming the fully-relativistic asymptotically matched operator from
$y=p/\delta$ and its historical collision-time normalization to the common $(p,\tau)$
variables of Sec.~\ref{sec:normalization} gives
exactly the same conservative radial/pitch form used by the finite-volume solver.  Using
the reference collision-frequency definition of Stahl et al.,
\begin{equation}
 \widetilde\nu_{ee}
 =\frac{16\sqrt{\pi}\,n_e e^4\ln\Lambda_0}
 {3m_e^2v_e^3(4\pi\epsilon_0)^2},
 \qquad
 (\tau_{\rm ref}\widetilde\nu_{ee})\frac{3\sqrt{\pi}}{4}
 =\chi_c\delta^{-3}.
 \label{eq:matched_numeric_frequency_conversion}
\end{equation}
The resulting thermal-log common-normalization base coefficients are
\begin{align}
 C_F^{\rm match,0}(p)
 &=\chi_c\,\frac{2\Psi(x)}{\delta^2},
 \label{eq:matched_numeric_cf}\\
 C_A^{\rm match,0}(p)
 &=\chi_c\,\frac{\gamma}{p}\Psi(x),
 \label{eq:matched_numeric_ca}\\
 \nu_{D,ee}^{\rm match,0}(p)
 &=\chi_c\,\frac{\gamma}{p^3}
 \left[
  \Phi(x)-\Psi(x)+\frac{\delta^2p^2}{2\gamma^2}
 \right],\\
 \nu_{D,ei}^{\rm match,0}(p)
 &=\chi_c\,\frac{\gamma}{p^3}Z_{\rm eff},\\
 \nu_D^{\rm match,0}(p)
 &=\nu_{D,ee}^{\rm match,0}(p)+\nu_{D,ei}^{\rm match,0}(p).
 \label{eq:matched_numeric_nud}
\end{align}
Here
\begin{equation}
 \Phi(x)=\operatorname{erf}(x),\qquad
 \Psi(x)=\frac{\Phi(x)-2x e^{-x^2}/\sqrt{\pi}}{2x^2}.
 \label{eq:matched_numeric_phi_psi}
\end{equation}

\paragraph{Common energy-dependent Coulomb-logarithm layer in solver normalization.}
At every radial evaluation point, compute the same physical ratios as in
Eq.~\eqref{eq:physical_log_ratios}, with $u=p$,
\begin{equation}
 R_{ee}(p)=\frac{\ln\Lambda_{ee}(p)}{\ln\Lambda_0},
 \qquad
 R_{ei}(p)=\frac{\ln\Lambda_{ei}(p)}{\ln\Lambda_0}.
 \label{eq:numeric_log_ratios}
\end{equation}
For either $b\in\{\mathrm{FT},\mathrm{match}\}$ the Coulomb-logarithm-corrected
complete-screening arrays are
\begin{align}
 C_F^{b,\rm CL}(p)
 &=R_{ee}(p)\,C_F^{b,0}(p),
 \label{eq:common_log_numeric_cf}\\
 C_A^{b,\rm CL}(p)
 &=C_A^{b,0}(p),
 \label{eq:common_log_numeric_ca}\\
 \nu_D^{b,\rm CL}(p)
 &=R_{ee}(p)\,\nu_{D,ee}^{b,0}(p)
  +R_{ei}(p)\,\nu_{D,ei}^{b,0}(p).
 \label{eq:common_log_numeric_nud}
\end{align}
This step is identical for the two base choices.  In particular, changing the small-angle
base model does not disable, replace, or redefine either energy-dependent Coulomb
logarithm, and $C_A$ is unchanged by this layer.

\paragraph{Common partial-screening overlay in solver normalization.}
The same screening layer is applied after the common energy-dependent Coulomb-logarithm
correction for either base coefficient set.  In the common normalized variables its
contribution is
\begin{align}
 \Delta C_F^{\rm PS}(p)
 &=\chi_c\,\frac{\gamma^2}{p^2}\frac{h(p)}{\ln\Lambda_0},
 \label{eq:screening_numeric_cf}\\
 \Delta C_A^{\rm PS}(p)&=0,
 \label{eq:screening_numeric_ca}\\
 \Delta\nu_D^{\rm PS}(p)
 &=\chi_c\,\frac{\gamma}{p^3}\frac{g(p)}{\ln\Lambda_0}.
 \label{eq:screening_numeric_nud}
\end{align}
Thus, for either $b\in\{\mathrm{FT},\mathrm{match}\}$,
\begin{equation}
 \boxed{
 C_F=C_F^{b,\rm CL}+\Delta C_F^{\rm PS},\qquad
 C_A=C_A^{b,\rm CL}=C_A^{b,0},\qquad
 \nu_D=\nu_D^{b,\rm CL}+\Delta\nu_D^{\rm PS}.
 }
 \label{eq:common_screened_coefficients}
\end{equation}
The two base evaluations above therefore feed the finite-volume assembly through exactly
the same Coulomb-logarithm and screening semantics.  When every ion is fully stripped,
$h=g=0$ and Eq.~\eqref{eq:common_screened_coefficients} reduces to the selected base with
the common energy-dependent Coulomb-logarithm correction still active.  If one also sets
$R_{ee}=R_{ei}=1$, the corresponding thermal-log base is recovered exactly.  Disabling
radial energy diffusion remains a separate explicit model switch and is applied after
coefficient construction.
Equations~\eqref{eq:matched_numeric_cf}--\eqref{eq:matched_numeric_nud} are a coordinate
transformation of the Section-1 fully-relativistic asymptotically matched thermal-log base,
not a new interpolation or fit.  The common logarithm corrections in
Eqs.~\eqref{eq:common_log_numeric_cf}--\eqref{eq:common_log_numeric_nud} and the common
partial-screening terms in Eqs.~\eqref{eq:screening_numeric_cf}--\eqref{eq:screening_numeric_nud}
are then applied independently of the selected base.  The local finite-volume assembly
receives the same three final coefficient arrays in either case.  The production default
keeps radial energy diffusion.  A benchmark configuration may deliberately set $C_A=0$
after coefficient evaluation; this is an explicit model switch and is never activated
implicitly by mesh, temperature, screening, or large-angle settings.

No rapidity quadrature or modified Bessel function is required for this configuration.
The only low-momentum numerical cancellation occurs in the direct expression for
$\Psi(x)$.  A fixed small-$x$ branch evaluates
\begin{equation}
 \Psi^{\rm ser}(x)
 =\frac{2x}{\sqrt{\pi}}
 \left(
  \frac13-\frac{x^2}{5}+\frac{x^4}{14}-\frac{x^6}{54}
  +\frac{x^8}{264}-\frac{x^{10}}{1560}
 \right),
 \label{eq:matched_psi_smallx}
\end{equation}
with the polynomial evaluated in $x^2$ by Horner's rule; the direct and series branches
must agree in an overlap interval.  For the production domain $p_{\min}=0$, no coefficient is required at the origin itself
because the inner radial flux is imposed directly.  On a shifted domain all coefficient
faces lie at finite momentum.  Thus there is no separate origin quadrature or singular
coefficient evaluation to perform.

For qualification, the transformed thermal-log base coefficients must reproduce the
Section-1 fully-relativistic asymptotically matched operator pointwise and recover its
intended asymptotic limits.  In particular, for $p\gg1$ with $\delta\ll1$,
\begin{equation}
 C_F^{\rm match,0}\rightarrow\chi_c,\qquad
 C_A^{\rm match,0}\rightarrow\frac{\chi_c\delta^2}{2},\qquad
 \nu_D^{\rm match,0}\rightarrow
 \chi_c\frac{1+Z_{\rm eff}}{p^2}
 \quad\text{to leading order in }\delta.
 \label{eq:matched_relativistic_limits}
\end{equation}
After the common Coulomb-logarithm layer, the corresponding complete-screening leading
behavior is instead
\begin{equation}
 C_F^{\rm match,CL}\rightarrow\chi_c R_{ee}(p),\qquad
 \nu_D^{\rm match,CL}\rightarrow
 \frac{\chi_c}{p^2}\left[R_{ee}(p)+Z_{\rm eff}R_{ei}(p)\right],
 \label{eq:matched_relativistic_limits_logs}
\end{equation}
while $C_A$ retains the base limit.  Partial screening adds the independent $h$ and $g$
terms afterward.  The non-relativistic thermal/suprathermal limit of the base is checked
directly against Eq.~\eqref{eq:matched_tp_operator} rather than by introducing a second
numerical formula.

\subsection{Conservative local flux discretization}

Integrating the differential part of Eq.~\eqref{eq:normalized_forward} over a cell gives
\begin{equation}
 \frac{dN_{ij}}{d\tau}\bigg|_{\rm local}
 =-\left(F^p_{i+1/2,j}-F^p_{i-1/2,j}\right)
  -\left(F^\xi_{i,j+1/2}-F^\xi_{i,j-1/2}\right).
 \label{eq:fv_balance}
\end{equation}
Each interior face flux is formed once and enters the two adjacent cells with opposite
signs.  Interior conservation is therefore algebraic rather than approximate.

The drift terms are discretized with first-order upwinding and the diffusion terms with
centered finite differences.  This deliberately simple choice is the baseline numerical
scheme for the present solver.  Higher-order or exponentially fitted Fokker--Planck
schemes are outside the scope of the present model and may be considered only after the
first-order formulation is fully qualified.

For an interior radial face, define
\begin{equation}
 A^p_{i+1/2,j}=A_p(p_{i+1/2},\xi_j),\qquad
 C_{A,i+1/2}=C_A^{({\rm sel})}(p_{i+1/2}),
\end{equation}
and the upwind face value
\begin{equation}
 f^{p,\mathrm{up}}_{i+1/2,j}=
 \begin{cases}
   f_{ij}, & A^p_{i+1/2,j}\ge 0,\\
   f_{i+1,j}, & A^p_{i+1/2,j}<0.
 \end{cases}
 \label{eq:radial_upwind_state}
\end{equation}
The integrated radial face flux is then
\begin{equation}
 F^p_{i+1/2,j}
 =2\pi p_{i+1/2}^2\Delta\xi
 \left[
 A^p_{i+1/2,j}f^{p,\mathrm{up}}_{i+1/2,j}
 -C_{A,i+1/2}
 \frac{f_{i+1,j}-f_{ij}}{\Delta p}
 \right].
 \label{eq:radial_fv_flux}
\end{equation}
Thus radial advection is first-order upwind, while physical energy diffusion is centered.

For an interior pitch face, define
\begin{equation}
 D_\xi(p,\xi)=\frac{\nu_D^{({\rm sel})}(p)}{2}(1-\xi^2),
\end{equation}
\begin{equation}
 A^\xi_{i,j+1/2}=A_\xi(p_i,\xi_{j+1/2}),\qquad
 D^\xi_{i,j+1/2}=D_\xi(p_i,\xi_{j+1/2}),
\end{equation}
and
\begin{equation}
 f^{\xi,\mathrm{up}}_{i,j+1/2}=
 \begin{cases}
   f_{ij}, & A^\xi_{i,j+1/2}\ge 0,\\
   f_{i,j+1}, & A^\xi_{i,j+1/2}<0.
 \end{cases}
 \label{eq:pitch_upwind_state}
\end{equation}
The integrated pitch face flux is
\begin{equation}
 F^\xi_{i,j+1/2}
 =\frac{2\pi}{3}
  \left(p_{i+1/2}^3-p_{i-1/2}^3\right)
 \left[
 A^\xi_{i,j+1/2}f^{\xi,\mathrm{up}}_{i,j+1/2}
 -D^\xi_{i,j+1/2}
 \frac{f_{i,j+1}-f_{ij}}{\Delta\xi}
 \right].
 \label{eq:pitch_fv_flux}
\end{equation}
The same first-order-upwind/centered-diffusion rule is therefore used in both momentum
coordinates.  Accuracy is obtained by mesh refinement rather than by introducing a more
elaborate spatial discretization at this stage.
The superscript ``sel'' emphasizes that the stencil is independent of the selected
small-angle collision model; only the precomputed coefficient values differ.

\subsection{Boundary conditions}

At the inner radial face two explicit configurations are available.  For the production
domain $p_{\min}=0$, regularity and the vanishing radial surface area impose
\begin{equation}
 F^p_{1/2,j}=0.
 \label{eq:inner_zero_flux}
\end{equation}
For a shifted domain $p_{\min}>0$, an absorbing zero-inflow/open-outflow condition may
instead be selected.  The exterior distribution is zero: advective inflow is suppressed,
advective drift with $A_p(p_{\min},\xi_j)<0$ leaves the represented domain, and, when
radial diffusion is enabled, the boundary value $f=0$ is imposed at the face using the
center-to-face distance $\Delta p/2$.  This is the boundary used by suprathermal
benchmark configurations.

At $\xi=\pm1$, both the pitch drift and diffusion coefficients vanish through
$1-\xi^2$, so
\begin{equation}
 F^\xi_{i,1/2}=F^\xi_{i,N_\xi+1/2}=0.
 \label{eq:pitch_boundary_flux}
\end{equation}

The outer radial boundary is a zero-inflow, open-outflow boundary.  The first-order
boundary flux is
\begin{equation}
 F^p_{N_p+1/2,j}
 =2\pi p_{\max}^2\Delta\xi\,
  \left[A_p(p_{\max},\xi_j)\right]_+ f_{N_p,j},
 \qquad [a]_+\equiv\max(a,0),
 \label{eq:outer_flux}
\end{equation}
with zero diffusive flux through the outer face and no incoming distribution supplied
from $p>p_{\max}$.  The value of $p_{\max}$ is chosen sufficiently large that the
reported solution is insensitive to further domain extension.  Domain-size convergence
is tested by increasing $p_{\max}$ while holding $\Delta p$ fixed.

\subsection{Discrete Chiu--Harvey large-angle gain}

The large-angle configuration is selected independently of the small-angle collision
model.  It is either disabled, $L_{\rm LA}=0$, or the Chiu--Harvey gain is enabled.  For
constant prescribed plasma parameters the Chiu--Harvey source is a fixed linear map of
the evolving finite-volume state and is retained in the same fully implicit BE/BDF2
operator.

\paragraph{Chiu--Harvey finite-energy aligned-primary gain.}
The physical Chiu--Harvey gain operator is fully defined in
Sec.~1, in particular by Eqs.~\eqref{eq:ch_primary_spectrum},
\eqref{eq:ch_primary_root}, and \eqref{eq:ch_reduced_source}, together with the
knock-on kinematic and cutoff conditions stated there.  The numerical model does not
introduce a different large-angle collision model; it only evaluates that operator on
the finite-volume mesh.

For each radial cell, the pitch-integrated primary distribution appearing in
Eq.~\eqref{eq:ch_primary_spectrum} is evaluated directly from the current 2P
solution as
\begin{equation}
 \boxed{
 \mathcal F_i(\tau)
 =\Delta\xi\sum_{j=1}^{N_\xi} f_{ij}(\tau)
 }
 \label{eq:discrete_primary_spectrum}
\end{equation}
which is the pitch integral of the piecewise-constant reconstruction in radial cell
$i$.  Thus the primary distribution used by the Chiu--Harvey source is the
pitch-integrated distribution of the evolving kinetic solution itself; no separate
primary distribution or analytic closure is introduced.

For a target cell center $(p_i,\xi_j)$, evaluate the required primary momentum
$p'_0(p_i,\xi_j)$ from Eq.~\eqref{eq:ch_primary_root}.  If the root and the associated
secondary satisfy all of the knock-on kinematic and large-angle-cutoff conditions of
Sec.~1, and if $p'_0$ lies inside the represented radial domain, let $k$ be the radial
cell containing that momentum,
\begin{equation}
 p_{k-1/2}\le p'_0(p_i,\xi_j)<p_{k+1/2}.
 \label{eq:ch_primary_cell}
\end{equation}
The first-order piecewise-constant representation gives
\begin{equation}
 \mathcal F\!\left(p'_0,\tau\right)\approx\mathcal F_k(\tau).
 \label{eq:ch_piecewise_primary}
\end{equation}
The normalized Chiu--Harvey source at the target cell center is then obtained by direct
substitution of this value into Eq.~\eqref{eq:ch_reduced_source}, using the normalization
and event-rate scale of Sec.~\ref{sec:normalization}.  Denote that value by
\begin{equation}
 S_{{\rm CH},ij}(\tau)
 =S_{\rm CH}[f](p_i,\xi_j,\tau).
 \label{eq:ch_cell_center_source}
\end{equation}
If the Section-1 kinematic or cutoff conditions are not satisfied, then
$S_{{\rm CH},ij}=0$.

The corresponding finite-volume contribution is evaluated with the same
piecewise-constant, first-order cell representation used for the rest of the baseline
solver,
\begin{equation}
 \boxed{
 \left.\frac{dN_{ij}}{d\tau}\right|_{\rm CH}
 =V_{ij}S_{{\rm CH},ij}
 }
 \label{eq:discrete_ch_source}
\end{equation}
where $V_{ij}$ is given by Eq.~\eqref{eq:cellvolume}.  Therefore the discrete
large-angle calculation is simply
\begin{equation}
 f_{ij}
 \;\longrightarrow\;
 \mathcal F_i=\Delta\xi\sum_j f_{ij}
 \;\longrightarrow\;
 S_{{\rm CH},ij}
 \;\longrightarrow\;
 \left.dN_{ij}/d\tau\right|_{\rm CH}.
\end{equation}
No additional physical approximation is introduced in Sec.~2.  Spatial accuracy of this
first-order evaluation is established by refinement of the uniform $p$ and $\xi$ meshes.
For fixed physical parameters and a fixed mesh, Eqs.~\eqref{eq:discrete_primary_spectrum}--
\eqref{eq:discrete_ch_source} define a time-independent linear mapping of the evolving
cell contents.


\paragraph{Exact discrete pitch-reduction factorization.}
For the algebra used by the GPU implementation it is convenient to introduce the radial
cell phase volume
\begin{equation}
 V_i^{(p)}=\frac{2\pi}{3}
 \left(p_{i+1/2}^3-p_{i-1/2}^3\right)
\end{equation}
and the radial particle content
\begin{equation}
 \boxed{
 M_i(\tau)\equiv\sum_{j=1}^{N_\xi}N_{ij}(\tau)
 =V_i^{(p)}\mathcal F_i(\tau).
 }
 \label{eq:ch_radial_mass}
\end{equation}
Let $R$ denote the fixed pitch-reduction map
\begin{equation}
 \bm M=R\bm N,
 \qquad
 (R\bm N)_i=\sum_{j=1}^{N_\xi}N_{ij}.
 \label{eq:ch_reduction_matrix}
\end{equation}
For every active Chiu--Harvey target cell $q\equiv(i,j)$, let $k(q)$ denote the
primary radial cell selected by Eq.~\eqref{eq:ch_primary_cell}.  All target geometry,
knock-on factors, normalization factors, target-cell volume, and the factor
$1/V_{k(q)}^{(p)}$ are time independent and may therefore be collected into a scalar
$c_q$.  The discrete source then has the exact form
\begin{equation}
 \left.\frac{dN_q}{d\tau}\right|_{\rm CH}
 =c_q M_{k(q)}.
 \label{eq:ch_mass_scatter}
\end{equation}
Define the fixed scatter map $U$ by $(U\bm M)_q=c_qM_{k(q)}$ on active targets and
zero otherwise.  The complete discrete Chiu--Harvey operator is therefore
\begin{equation}
 \boxed{L_{\rm CH}=UR.}
 \label{eq:ch_UR_factorization}
\end{equation}
Equations~\eqref{eq:ch_radial_mass}--\eqref{eq:ch_UR_factorization} are not a new
quadrature or a new physical approximation: they are an exact factorization of
Eqs.~\eqref{eq:discrete_primary_spectrum}--\eqref{eq:discrete_ch_source}.  In particular,
the source can be applied by one pitch reduction per radial shell followed by one scatter
to each active target instead of repeating the same pitch sum separately for every target.
The expanded and reduced source applications are required to agree to floating-point
accuracy.

\subsection{Semidiscrete system and discrete balance}

Flatten the one-based cell indices with
\begin{equation}
 k(i,j)=(i-1)N_\xi+j,
 \qquad k=1,\ldots,N_pN_\xi.
 \label{eq:flattening}
\end{equation}
After discretizing every local differential term, let $L_0$ denote the time-independent
finite-volume operator containing electric acceleration, the selected small-angle
collision coefficients, synchrotron reaction, and all boundary fluxes.  The selected
large-angle operator is
\begin{equation}
 L_{\rm LA}=\begin{cases}
 0, & \text{no large-angle gain},\\
 U_{\rm CH}R_{\rm CH}, & \text{Chiu--Harvey}.
 \end{cases}
 \label{eq:large_angle_discrete_choices}
\end{equation}
The complete semidiscrete equation is
\begin{equation}
 \boxed{
 L=L_0+L_{\rm LA},\qquad
 \frac{d\bm N}{d\tau}=L\bm N.
 }
 \label{eq:semidiscrete}
\end{equation}
This is one fixed linear operator; no operator splitting is introduced.

Summing the physical finite-volume equations cancels every shared interior conservative
flux exactly.  Writing $\Phi_{\rm in}\ge0$ and $\Phi_{\rm out}\ge0$ for particle loss
through the inner and outer radial boundaries, respectively, gives
\begin{equation}
 \boxed{
 \frac{d}{d\tau}\bm 1^T\bm N
 =-\Phi_{\rm in}-\Phi_{\rm out}+\bm 1^TL_{\rm LA}\bm N.
 }
 \label{eq:discrete_particle_balance}
\end{equation}
For the zero-flux production inner boundary $\Phi_{\rm in}=0$; it can be nonzero for the
shifted absorbing boundary.  For Chiu--Harvey the production term is
$\bm1^TU_{\rm CH}R_{\rm CH}\bm N$, and it vanishes when the gain is disabled.

\subsection{Fully implicit time stepping}

The normalized time grid is uniform,
\begin{equation}
 \tau^n=n\Delta\tau,\qquad n=0,1,\ldots,N_t.
 \label{eq:uniform_time_grid}
\end{equation}
The complete selected semidiscrete operator of Eq.~\eqref{eq:semidiscrete} is advanced
implicitly.  No enabled physical term is lagged, split off, or treated explicitly.
A single backward-Euler step initializes the two-step method,
\begin{equation}
 \boxed{\left(I-\Delta\tau L\right)\bm N^1=\bm N^0}
 \label{eq:be}
\end{equation}
and every subsequent step uses BDF2,
\begin{equation}
 \boxed{
 \left(\frac{3}{2}I-\Delta\tau L\right)\bm N^{n+1}
 =2\bm N^n-\frac{1}{2}\bm N^{n-1},\qquad n\ge1.
 }
 \label{eq:bdf2}
\end{equation}
For fixed physical parameters, grid, and timestep, the BE and BDF2 matrices are constant
and differ only in the coefficient multiplying the physical identity.

When large-angle gain is disabled, Eqs.~\eqref{eq:be}--\eqref{eq:bdf2} are solved
directly on the physical state.  With Chiu--Harvey enabled, introduce the radial
auxiliary state $\bm M=R_{\rm CH}\bm N$.  The exact augmented systems are
\begin{equation}
 \boxed{
 \begin{bmatrix}
  I-\Delta\tau L_0 & -\Delta\tau U_{\rm CH}\\
  -R_{\rm CH} & I
 \end{bmatrix}
 \begin{bmatrix}\bm N^1\\\bm M^1\end{bmatrix}
 =\begin{bmatrix}\bm N^0\\\bm0\end{bmatrix}
 }
 \label{eq:be_augmented}
\end{equation}
and
\begin{equation}
 \boxed{
 \begin{bmatrix}
  \frac32 I-\Delta\tau L_0 & -\Delta\tau U_{\rm CH}\\
  -R_{\rm CH} & I
 \end{bmatrix}
 \begin{bmatrix}\bm N^{n+1}\\\bm M^{n+1}\end{bmatrix}
 =\begin{bmatrix}2\bm N^n-\frac12\bm N^{n-1}\\\bm0\end{bmatrix}.
 }
 \label{eq:bdf2_augmented}
\end{equation}
Eliminating $\bm M$ recovers the physical systems with
$L=L_0+U_{\rm CH}R_{\rm CH}$ exactly.  The auxiliary constraint contains no time
derivative, so its identity block is unchanged between BE and BDF2.

Implicit treatment removes explicit stability restrictions associated with stiff
collision and advection terms, but not the timestep requirement for accuracy.  The
accepted $\Delta\tau$ is therefore established by temporal refinement, and cell contents
are monitored for loss of non-negativity.

\subsection{Sparse direct solution on the GPU}

The local physical block always has the same finite-volume stencil.  Changing between
the two small-angle collision models changes coefficient values but not the CSR topology.
The Chiu--Harvey configuration adds the exact low-dimensional radial auxiliary system of
Eqs.~\eqref{eq:be_augmented}--\eqref{eq:bdf2_augmented}.

For one constant-parameter simulation, the direct-solve procedure is:
\begin{enumerate}
 \item construct the uniform finite-volume grid and cell geometry;
 \item project the prescribed initial distribution using Eq.~\eqref{eq:initial_projection};
 \item evaluate the selected small-angle collision coefficients once;
 \item if Chiu--Harvey is enabled, construct its fixed target geometry and the exact
       reduction/scatter maps $R_{\rm CH}$ and $U_{\rm CH}$;
 \item construct and validate the fixed CSR graph for the physical or augmented system;
 \item assemble the backward-Euler matrix;
 \item perform one sparse symbolic analysis of that fixed graph;
 \item numerically factor the backward-Euler matrix and solve the startup step;
 \item change only the physical diagonal coefficient from $1$ to $3/2$;
 \item numerically factor the BDF2 matrix once while reusing symbolic analysis, then use
       the retained factorization for all subsequent timesteps.
\end{enumerate}
After startup, every timestep requires only the BDF2 right-hand side and one sparse
triangular solve with the retained factorization.

For Chiu--Harvey, the exact augmented graph contains the local physical stencil, one
auxiliary-column coupling for each active target, and $N_p$ radial constraints
$M_i-\sum_jN_{ij}=0$.  Its raw graph size is $O(N_pN_\xi)$ and avoids the quadratic
pitch growth that would result from explicitly multiplying $U_{\rm CH}R_{\rm CH}$.

The implementation uses Warp-owned FP64 CSR matrices and state vectors on the CUDA
device.  CSR column indices, finite-volume slot maps, auxiliary couplings, and graph
validation are constructed with Warp GPU kernels; only the $O(N)$ CSR row-pointer scan
and small time-independent preprocessing remain on the host.  Raw CUDA pointers are
passed directly to cuDSS on the same CUDA stream.  cuDSS performs sparse symbolic
analysis, numerical factorization, and triangular solves.  No CPU kinetic matrix or CPU
sparse solver participates in production time integration.

For any solved system $A\bm x=\bm b$, the relative residual
\begin{equation}
 r_{\rm lin}=\frac{\|A\bm x-\bm b\|_2}{\|\bm b\|_2}
 \label{eq:linear_residual}
\end{equation}
is evaluated directly.  When the Chiu--Harvey augmented formulation is active, the
constraint defect $\|\bm M-R_{\rm CH}\bm N\|$ is checked independently.

\subsection{Discrete diagnostics and numerical qualification}

The solver reports diagnostics that are direct moments of the represented kinetic
distribution.  Neither the no-gain nor Chiu--Harvey configuration requires a
threshold-based runaway definition.

The normalized represented density is exactly
\begin{equation}
 \boxed{
 \frac{n_F}{n_{\rm ref}}=\sum_{i=1}^{N_p}\sum_{j=1}^{N_\xi}N_{ij}
 }
 \label{eq:discrete_density}
\end{equation}
because $N_{ij}$ is the particle content in the normalized phase-space measure.

The dimensional parallel current density corresponding to the piecewise-constant
reconstruction is
\begin{equation}
 \boxed{
 j_\parallel
 =-e n_{\rm ref}c
 \sum_{i=1}^{N_p}\sum_{j=1}^{N_\xi}
 f_{ij}\,W_i^{(v)}W_j^{(\xi)}
 }
 \label{eq:discrete_current}
\end{equation}
with exact cell weights
\begin{align}
 W_i^{(v)}
 &=2\pi\int_{p_{i-1/2}}^{p_{i+1/2}}
       p^2\frac{p}{\sqrt{1+p^2}}\,dp
 \nonumber\\
 &=2\pi\left[
       \frac{\gamma^3}{3}-\gamma
       \right]_{p_{i-1/2}}^{p_{i+1/2}},
 \qquad \gamma=\sqrt{1+p^2},
 \label{eq:current_radial_weight}\\
 W_j^{(\xi)}
 &=\int_{\xi_{j-1/2}}^{\xi_{j+1/2}}\xi\,d\xi
 =\frac{\xi_{j+1/2}^2-\xi_{j-1/2}^2}{2}.
 \label{eq:current_pitch_weight}
\end{align}
Thus neither density nor current requires an additional diagnostic quadrature.  Other
moments may be added later by integrating their physical weight against the same
piecewise-constant reconstruction, but they are not part of the baseline numerical
specification.

\paragraph{Numerical refinement parameters.}
Every selected physical configuration is qualified by independent refinement of
\begin{equation}
 \boxed{
 \Delta p,\qquad \Delta\xi,\qquad p_{\max},\qquad
 q_{\rm init},\qquad \Delta\tau
 }
 \label{eq:convergence_parameters}
\end{equation}
with $q_{\rm coll}$ added as an independent convergence parameter only for the
finite-temperature relativistic coefficient evaluation.  The fully-relativistic asymptotically matched coefficients are analytic apart
from their fixed small-$x$ series branch and therefore
use no collision quadrature.  Fixed series orders and branch switches are implementation
choices, not user convergence parameters.

Spatial refinement is performed on the uniform grid family of
Eqs.~\eqref{eq:uniform_p_grid}--\eqref{eq:uniform_xi_grid}.  Radial and pitch resolution
are tested independently by reducing $\Delta p$ at fixed $\Delta\xi$ and reducing
$\Delta\xi$ at fixed $\Delta p$.  Domain-size convergence is a separate test:
$p_{\max}$ is increased while $\Delta p$ is held fixed, so that increasing the physical
domain does not simultaneously coarsen the radial mesh.  The accepted $p_{\max}$ is
large enough that the represented distribution and reported moments are insensitive to
a further increase of the outer boundary over the simulated time interval.

For the finite-temperature model, $q_{\rm coll}$ is increased until the collision
coefficients are unchanged to the accuracy required by the mesh-converged kinetic
solution.  The fully-relativistic asymptotically matched model instead requires overlap and
asymptotic checks of its analytic direct/series evaluation.  In all configurations, $q_{\rm init}$ is increased
until the projected initial cell contents in Eq.~\eqref{eq:initial_projection} are
converged.  These checks are performed independently of mesh refinement.

\paragraph{Required numerical qualification tests.}
These tests establish that each implemented configuration discretizes its Section-1
operator correctly.  They are distinct from external literature benchmarks, which test
whether the resulting physical model reproduces established kinetic results.

\begin{enumerate}
 \item \textbf{Small-angle coefficient evaluation.}
       For the finite-temperature model, the direct scaled formulas and low-$p$ Taylor
       formulas must agree in an overlap interval and increasing $q_{\rm coll}$ must
       leave the coefficients unchanged.  For the fully-relativistic asymptotically matched model,
       Eqs.~\eqref{eq:matched_numeric_cf}--\eqref{eq:matched_numeric_nud} must reproduce the
       Section-1 operator in Eq.~\eqref{eq:matched_tp_operator} after coordinate
       transformation; the direct and small-$x$ evaluations of $\Psi$ must overlap, and
       the non-relativistic and relativistic asymptotic limits must be recovered.
       For both base models, the energy-dependent Coulomb-logarithm layer must satisfy
       Eqs.~\eqref{eq:common_log_numeric_cf}--\eqref{eq:common_log_numeric_nud}; setting
       $R_{ee}=R_{ei}=1$ must recover the thermal-log base exactly, and $C_A$ must remain
       unchanged.  Fully stripped ions must set only the partial-screening terms to zero,
       leaving the common Coulomb-logarithm corrections active.  With bound electrons
       present, the additional changes in $C_F$ and $\nu_D$ must equal
       Eqs.~\eqref{eq:screening_numeric_cf} and \eqref{eq:screening_numeric_nud},
       respectively, again with $C_A$ unchanged.

 \item \textbf{Finite-volume balance.}
       For arbitrary test states, summing the semidiscrete physical-cell equations must
       reproduce Eq.~\eqref{eq:discrete_particle_balance} to floating-point accuracy.
       Every interior face contribution must cancel pairwise between adjacent cells.

 \item \textbf{Collision equilibrium.}
       With electric field, synchrotron radiation, and large-angle gain removed, the
       appropriate Maxwell--J{\"u}ttner equilibrium at the prescribed background
       temperature must approach a stationary discrete solution as the uniform mesh and
       any applicable coefficient quadrature are refined.  Exact preservation is not
       required from the first-order spatial scheme; the equilibrium defect must converge
       to zero.

 \item \textbf{Spatial convergence of the kinetic equation.}
       Smooth verification problems must recover first-order asymptotic convergence with
       respect to the upwinded directions as $\Delta p$ and $\Delta\xi$ are independently
       reduced.  Production configurations must then be continued to finer meshes until
       the full distribution and reported moments are insensitive to further refinement.

 \item \textbf{Chiu--Harvey discretization, when selected.}
       For a prescribed smooth two-dimensional distribution, the discrete primary pitch
       integral and the resulting cell-centered source must converge to the Section-1
       Chiu--Harvey expressions under independent $\Delta p$ and $\Delta\xi$ refinement.
       The test includes the finite-energy primary root, kinematic support, and fixed
       large-angle cutoff.

 \item \textbf{Chiu--Harvey algebraic factorization, when selected.}
       On a fixed mesh and state, the pitch-reduction/scatter application must agree with
       the explicitly expanded $N_\xi$-term source sum to floating-point accuracy.  The
       augmented BE and BDF2 systems must reproduce the physical solution of the expanded
       systems to the accuracy permitted by the sparse residual, and
       $\bm M=R_{\rm CH}\bm N$ must hold to floating-point accuracy.

 \item \textbf{Outer-domain convergence.}
       At fixed $\Delta p$ and $\Delta\xi$, increasing $p_{\max}$ must leave the solution
       on the original momentum interval and the reported moments unchanged to the
       required accuracy.

 \item \textbf{Temporal convergence.}
       With a spatially converged mesh, reducing $\Delta\tau$ must recover second-order
       convergence of the fully implicit BDF2 evolution over a fixed physical interval,
       including the single backward-Euler startup step.

 \item \textbf{Sparse direct solve.}
       Every BE and BDF2 solve must satisfy Eq.~\eqref{eq:linear_residual}; the residual
       must remain small compared with the demonstrated spatial and temporal
       discretization errors.  Any active auxiliary constraint must satisfy its separate
       defect check.
\end{enumerate}

Unresolved coarse meshes may be used for algebraic, indexing, and implementation tests,
but they are not evidence of physical convergence.  The numerical strategy is deliberately to use a simple first-order finite-volume
discretization on sufficiently large uniform GPU meshes and to demonstrate convergence
directly for each selected physical configuration.


%===============================================================================
\section{Future bulk-plasma coupling}
\label{sec:future_bulk_coupling}

The standalone solver currently advances the kinetic distribution using a prescribed,
constant plasma state.  A future coupled driver will evolve a bulk plasma model and
exchange the quantities needed to update the kinetic coefficients and source terms.
That extension must be defined explicitly rather than inferred from the present
configuration file.

At minimum, the coupled formulation must specify:
\begin{itemize}
  \item the bulk state variables exchanged with the kinetic solver, including density,
        temperature, electric field, magnetic field, ionization state, and any source
        or loss rates;
  \item the coupling cadence and whether the bulk state is held fixed during a kinetic
        step or iterated to a coupled time-level convergence;
  \item the coefficient and Chiu--Harvey source update, including which GPU arrays are
        updated in place and which require host preprocessing;
  \item conservation equations for particle, energy, and charge exchange, with a
        defined sign convention and diagnostics on both sides of the interface;
  \item the sparse-topology policy: unchanged mesh topology may reuse cuDSS analysis
        and use numerical refactorization, while changed sparsity requires rebuilding
        and requalifying the CSR graph; and
  \item restart data containing the bulk state, kinetic state, time level, configuration,
        solver metadata, and enough provenance to reproduce the coupling decision.
\end{itemize}

The coupled implementation should first be qualified against a prescribed-state run
using identical coefficients.  It should then demonstrate conservation and coefficient
parity over representative bulk states, compare topology reuse with full rebuilds, and
test timestep, mesh, and coupling-cadence refinement.  Until those checks are complete,
the standalone forward solver remains the reference calculation.

\begin{thebibliography}{99}

\bibitem{ConnorHastie1975}
J.~W. Connor and R.~J. Hastie,
``Relativistic limitations on runaway electrons,''
\emph{Nucl. Fusion} \textbf{15}, 415 (1975).

\bibitem{BraamsKarney1989}
B.~J. Braams and C.~F.~F. Karney,
``Conductivity of a relativistic plasma,''
\emph{Phys. Fluids B} \textbf{1}, 1355 (1989).

\bibitem{PikeRose2014}
O.~J. Pike and S.~J. Rose,
``Dynamical friction in a relativistic plasma,''
\emph{Phys. Rev. E} \textbf{89}, 053107 (2014).

\bibitem{Hesslow2017Screening}
L.~Hesslow, O.~Embr\'eus, A.~Stahl, T.~C. DuBois, G.~Papp,
S.~L. Newton, and T.~F{\"u}l{\"o}p,
``Effect of partially screened nuclei on fast-electron dynamics,''
\emph{Phys. Rev. Lett.} \textbf{118}, 255001 (2017).

\bibitem{Hesslow2018Collision}
L.~Hesslow, O.~Embr\'eus, M.~Hoppe, T.~C. DuBois, G.~Papp,
M.~Rahm, and T.~F{\"u}l{\"o}p,
``Generalized collision operator for fast electrons interacting with partially ionized impurities,''
\emph{J. Plasma Phys.} \textbf{84}, 905840605 (2018).

\bibitem{Andersson2001}
F.~Andersson, P.~Helander, and L.-G.~Eriksson,
``Damping of relativistic electron beams by synchrotron radiation,''
\emph{Phys. Plasmas} \textbf{8}, 5221 (2001).

\bibitem{Stahl2016}
A.~Stahl, O.~Embr{\'e}us, G.~Papp, M.~Landreman, and T.~F{\"u}l{\"o}p,
``Kinetic modelling of runaway electrons in dynamic scenarios,''
\emph{Nucl. Fusion} \textbf{56}, 112009 (2016).

\bibitem{Moller1932}
C.~M{\o}ller,
``Zur Theorie des Durchgangs schneller Elektronen durch Materie,''
\emph{Ann. Phys.} \textbf{406}, 531 (1932).

\bibitem{Chiu1998}
S.~C. Chiu, M.~N. Rosenbluth, R.~W. Harvey, and V.~S. Chan,
\emph{Nucl. Fusion} \textbf{38}, 1711 (1998).

\bibitem{Harvey2000}
R.~W. Harvey, V.~S. Chan, S.~C. Chiu, T.~E. Evans,
M.~N. Rosenbluth, and D.~G. Whyte,
\emph{Phys. Plasmas} \textbf{7}, 4590 (2000).


\end{thebibliography}

\end{document}
